\documentclass[11pt,a4paper]{article}

\usepackage[T1]{fontenc}
\usepackage[utf8]{inputenc}
\usepackage[french]{babel}
\usepackage{amsmath,amssymb,mathtools,geometry,booktabs,xcolor,tikz,hyperref,listings}

\usetikzlibrary{arrows.meta,positioning,shapes.geometric,calc}
\geometry{margin=22mm}
\definecolor{navy}{RGB}{18,49,88}
\hypersetup{colorlinks=true,linkcolor=navy}
\lstset{basicstyle=\ttfamily\small,breaklines=true,frame=single}
\title{Réseaux de neurones et solveurs numériques\\\large Concepts, méthodes et implémentations}
\author{Frédéric Magoulès}
\date{Novembre 2026}

\begin{document}\maketitle\tableofcontents\newpage

\section{Organisation des ressources et des travaux pratiques}

Le répertoire \texttt{reference/} contient les \textbf{corrigés exécutables} des exercices : les notebooks complets permettent de vérifier les implémentations et de reproduire les expériences numériques.
Le répertoire \texttt{student/} contient les versions distribuées aux participants, avec les portions à compléter.
Le répertoire \texttt{instructor/} rassemble les \textbf{explications de compréhension} : fondements mathématiques, interprétations numériques, schémas et conseils pédagogiques. Il complète les corrigés sans les remplacer.
\begin{center}\begin{tabular}{ll}\toprule Notebook & Objectif \\\midrule
\texttt{model\_architecture} & MLP, CNN, graphe kNN, EdgeConv et GNN\\
\texttt{training\_process} & contraintes faibles et fortes de divergence\\
\texttt{hybrid\_corrector} & apprentissage résidu-vers-erreur\\
\texttt{hybrid\_accelerator} & apprentissage second membre-vers-solution\\\bottomrule\end{tabular}\end{center}

\section{Apprentissage supervisé avec PyTorch}

\subsection{Principes et mise en œuvre de l'apprentissage}

On dispose d'un ensemble de $N$ exemples $(x_i,y_i)$, où $x_i$ désigne une donnée d'entrée et $y_i$ la valeur attendue. Le problème consiste à construire une fonction capable de prédire $y_i$ à partir de $x_i$. Le réseau de neurones $F_\theta$ réalise cette prédiction : $F$ désigne la fonction représentée par le réseau, et $\theta$ l'ensemble de ses paramètres ajustables (poids et biais). L'indice $i$ numérote les exemples et $N$ est leur nombre total.

On cherche les paramètres optimaux $\theta^*$ en minimisant la fonction de coût $\mathcal L$ (\textit{loss function}), par exemple l'erreur quadratique moyenne :
\[
\boxed{\theta^*=\arg\min_\theta\mathcal L(\theta),\qquad
\mathcal L(\theta)=\frac{1}{N}\sum_{i=1}^N\|F_\theta(x_i)-y_i\|_2^2.}
\]
La norme $\|\cdot\|_2$ mesure l'écart entre la prédiction et la cible. Une époque (\textit{epoch}) correspond à un passage complet sur les données d'entraînement. Les données sont regroupées en mini-lots (\textit{mini-batches}) ; pour chaque mini-lot (\textit{mini-batch}), on calcule les prédictions par propagation avant (\textit{forward pass}), la perte (\textit{loss}), les gradients par rétropropagation (\textit{backpropagation}), puis la mise à jour des paramètres par l'optimiseur (\textit{optimizer}). Le taux d'apprentissage (\textit{learning rate}) règle l'amplitude des mises à jour.
\begin{lstlisting}[language=Python]
model = NeuralNetwork()  # A user-defined torch.nn.Module
optimizer = torch.optim.Adam(model.parameters(), lr=1e-3)
loss_fn = torch.nn.MSELoss()
loader = DataLoader(TensorDataset(X_train, Y_train),
                    batch_size=32, shuffle=True)
for epoch in range(150):
    model.train()
    for xb, yb in loader:
        pred = model(xb)        # Forward pass
        loss = loss_fn(pred,yb) # Objective function
        optimizer.zero_grad()  # Reset accumulated gradients
        loss.backward()        # Automatic differentiation
        optimizer.step()       # Update parameters
\end{lstlisting}
\paragraph{À retenir.} \texttt{loss.backward()} calcule les gradients mais ne modifie pas les paramètres ; \texttt{optimizer.step()} réalise la mise à jour. Pour une descente de gradient élémentaire, $\theta^{k+1}=\theta^k-\eta\nabla_\theta\mathcal L_B(\theta^k)$ ; Adam utilise des moments adaptatifs. Une courbe de validation, calculée sur des données non utilisées pour l'entraînement, permet d'évaluer la capacité du réseau à prédire sur de nouveaux points et de détecter un éventuel surapprentissage.

\subsection{Exemple introductif : approximation d'une fonction sinus par un MLP (\textit{mlp\_sine\_intro})}

On considère une fonction scalaire simple :
\[y(x)=\sin(2\pi x),\qquad x\in[0,1].\]
Le réseau $F_\theta:\mathbb R\to\mathbb R$ comporte deux couches cachées de 32 neurones avec activation \texttt{Tanh}. On minimise l'erreur quadratique moyenne sur des couples $(x_i,y_i)$ :
\[\mathcal L(\theta)=\frac1N\sum_{i=1}^{N}|F_\theta(x_i)-\sin(2\pi x_i)|^2.\]
Le programme utilise Adam, des mini-lots (\textit{mini-batches}) de 32 exemples et 200 époques (\textit{epochs}). Il trace deux résultats : les pertes d'entraînement et de validation en fonction des époques (\textit{epochs}), puis la fonction exacte et la prédiction du MLP. Les données de validation sont distinctes de celles d'entraînement.
\paragraph{Exécution.} Dans \texttt{instructor/}, lancer \texttt{python mlp\_sine\_intro.py} (dépendances : NumPy, PyTorch, Matplotlib), ou exécuter \texttt{mlp\_sine\_intro.ipynb}. L'image \texttt{mlp\_sine\_results.png} est enregistrée automatiquement.
\paragraph{Lien avec la suite.} Cet exemple introduit exclusivement la mécanique \texttt{forward} $\rightarrow$ \texttt{loss} $\rightarrow$ \texttt{backward} $\rightarrow$ \texttt{step}. Le notebook \texttt{model\_architecture} compare les architectures ; \texttt{training\_process} introduit les contraintes physiques ; \texttt{hybrid\_corrector} et \texttt{hybrid\_accelerator} apprennent respectivement une erreur et une solution de systèmes linéaires.

\section{Architectures de réseaux de neurones : MLP, CNN et GNN (\textit{model\_architecture})}

Dans \texttt{model\_architecture}, le MLP travaille sur une représentation aplatie/ponctuelle, le CNN exploite des convolutions locales sur une grille ordonnée, et le GNN agrège des messages selon \texttt{edge\_index}. Le graphe est construit par \texttt{create\_knn\_graph} sur les coordonnées irrégulières, avec $k=5$ par défaut. 
Chaque arête du maillage est représentée dans les deux orientations
afin de permettre la propagation des informations dans les deux sens.
Pour chaque nœud $i$, les messages provenant de ses voisins sont
calculés puis agrégés selon :
\[
m_{ij} = \operatorname{MLP}_{\theta}([h_i,h_j]),
\qquad
h'_i = \frac{1}{|\mathcal{N}(i)|}
\sum_{j \in \mathcal{N}(i)} m_{ij},
\]
où :
\begin{itemize}
    \item $i$ désigne le nœud dont on souhaite actualiser les
    caractéristiques ;
    \item $j$ désigne un nœud voisin de $i$ ;
    \item $h_i$ et $h_j$ sont les vecteurs de caractéristiques
    (\textit{features}) respectifs des nœuds $i$ et $j$ ;
    \item $[h_i,h_j]$ désigne la concaténation de ces deux vecteurs ;
    \item $\operatorname{MLP}_{\theta}$ est un réseau de neurones
    multicouche dont les paramètres sont notés $\theta$ ;
    \item $m_{ij}$ est le message calculé à partir des
    caractéristiques des nœuds $i$ et $j$, destiné au nœud $i$ ;
    \item $\mathcal{N}(i)$ est l'ensemble des voisins du nœud $i$,
    et $|\mathcal{N}(i)|$ leur nombre ;
    \item $h'_i$ est le nouveau vecteur de caractéristiques du
    nœud $i$, obtenu en calculant la moyenne des messages reçus.
\end{itemize}

Il convient d'éviter les doublons d'arêtes orientées, car ils
conduiraient à comptabiliser plusieurs fois certains messages
et modifieraient ainsi leur poids relatif dans l'agrégation
par moyenne.

Cette opération constitue une étape de propagation de messages
(\textit{message passing}) dans un réseau de neurones sur
graphes (GNN).
Le GNN possède trois couches \texttt{EdgeConv}; hors caractéristiques d'entrée globalement informatives, trois passages ne propagent pas les messages au-delà d'un voisinage de trois arêtes. Ici, la distance à la source est déjà fournie comme caractéristique à chaque nœud : la portée effective de l'information disponible ne se résume donc pas au nombre de couches.
\begin{center}\begin{tikzpicture}[>=Latex, every node/.style={font=\small}]
\foreach \x/\lab in {0/0,1.6/1,3.2/2,4.8/3,6.4/4,8/5}{\node[circle,draw=navy,minimum size=7mm] (n\lab) at (\x,0) {\lab};}
\foreach \a/\b in {0/1,1/2,2/3,3/4,4/5}{\draw[thick,navy] (n\a)--(n\b);}
\draw[->,thick,red] (n0.north) to[bend left=30] node[above] {3 couches : 3 sauts au plus} (n3.north);
\end{tikzpicture}\end{center}

\section{Apprentissage et accélération des solveurs numériques}

\subsection{Formulation du problème discret}

Considérons une discrétisation du problème de Poisson, avec conditions de Dirichlet homogènes :
\[-\Delta u=f\quad\text{dans }\Omega,\qquad u=0\quad\text{sur }\partial\Omega.\]
Après élimination des inconnues de bord, on obtient un système linéaire inversible
\[Au=f,\qquad u^*=A^{-1}f.\]
L'écriture $A^{-1}$ est ici \emph{conceptuelle} : les notebooks calculent des solutions de référence au moyen d'un solveur direct creux; ils ne forment pas la matrice inverse. L'opérateur $A$ est fixé par le maillage, la discrétisation et les conditions aux limites. Les réseaux apprennent des approximations de son action sur une \emph{distribution} de vecteurs.

\subsection{Méthode de Jacobi et propriétés de lissage}

Écrivons $A=D-(L+U)$, où $D$ est la diagonale et $L+U=D-A$. Jacobi s'écrit
\[u^{k+1}=u^k+D^{-1}(f-Au^k),\qquad e^{k+1}=(I-D^{-1}A)e^k.\]
La matrice d'itération est $T_J=I-D^{-1}A$. Pour le Laplacien modèle sur grille régulière, Jacobi amortit rapidement certaines composantes oscillantes, tandis que les modes lisses peuvent décroître lentement. 

\subsection{Accélération par apprentissage de la solution (\textit{hybrid\_accelerator})}

Dans \texttt{hybrid\_accelerator}, les données sont des paires $(f,u^*)$ où $Au^*=f$. Le CNN $G_\theta$ cherche à approximer
\[\boxed{G_\theta(f)\simeq A^{-1}f=u^*.}\]
La sortie est utilisée \emph{une seule fois} comme initialisation : $u^0=G_\theta(f)$, puis Jacobi effectue les itérations classiques. Le CNN n'est pas appliqué à chaque itération.
\begin{center}\begin{tikzpicture}[node distance=12mm,>=Latex, every node/.style={font=\small},box/.style={draw=navy,rounded corners,minimum height=9mm,align=center}]
\node[box] (f) {$f$};\node[box,right=of f] (cnn) {CNN $G_\theta$\\approx. de $A^{-1}$};\node[box,right=of cnn] (u) {$u^0$};\node[box,right=of u] (jac) {Jacobi\\$u^{k+1}=u^k+D^{-1}r^k$};\node[box,right=of jac] (res) {$u^m$};
\draw[->,thick] (f)--(cnn);\draw[->,thick] (cnn)--(u);\draw[->,thick] (u)--(jac);\draw[->,thick] (jac)--(res);
\end{tikzpicture}\end{center}

\subsection{Correction par apprentissage de l'erreur (\textit{hybrid\_corrector})}

Pour une approximation $u^k$, définissons
\[r^k=f-Au^k,\qquad e^k=u^*-u^k.\]
Alors, $Ae^k=r^k$ et
\[\boxed{e^k=A^{-1}r^k,\qquad C_\theta(r^k)\simeq e^k.}\]
Le notebook \texttt{hybrid\_corrector} apprend des paires $(r^k,e^k)$ obtenues après un nombre aléatoire d'itérations de Jacobi. Après \texttt{cnn\_at\_step} itérations, il applique une seule correction additive
\[u^{k,+}=u^k+C_\theta(r^k),\qquad e^{k,+}=e^k-C_\theta(r^k),\]
puis reprend Jacobi. Une correction exacte donnerait $u^{k,+}=u^*$, mais le réseau est approximatif : le résidu et l'erreur ne s'annulent pas nécessairement.
\begin{center}\begin{tikzpicture}[node distance=10mm,>=Latex, every node/.style={font=\small},box/.style={draw=navy,rounded corners,minimum height=9mm,align=center}]
\node[box] (uk) {$u^k$};\node[box,right=of uk] (r) {$r^k=f-Au^k$};\node[box,right=of r] (cnn) {CNN $C_\theta$\\approx. de $A^{-1}$};\node[box,right=of cnn] (e) {$\widehat e^k$};\node[box,right=of e] (plus) {$u^{k,+}=u^k+\widehat e^k$};
\draw[->,thick] (uk)--(r);\draw[->,thick] (r)--(cnn);\draw[->,thick] (cnn)--(e);\draw[->,thick] (e)--(plus);\draw[->,dashed] (uk.south) -- ++(0,-.7) -| (plus.south);
\end{tikzpicture}\end{center}

\subsection{Normalisation des données d'entrée et de sortie}

Dans \texttt{hybrid\_accelerator}, $s_f=\sqrt{\operatorname{mean}(f_{\rm int}^2)}+10^{-10}$; on entraîne sur $(f/s_f,u^*/s_f)$. Dans \texttt{hybrid\_corrector}, $s_r=\sqrt{\operatorname{mean}(r_{\rm int}^2)}+10^{-10}$; on entraîne sur $(r/s_r,e/s_r)$. En inférence, on multiplie la prédiction par le même facteur. Cette construction respecte l'homogénéité de l'opérateur linéaire cible, sans garantir que le CNN appris soit lui-même linéaire.

\section{Apprentissage sous contraintes physiques (\textit{training\_process})}

Le notebook source s'appelle \texttt{training\_process.ipynb}. Il reconstruit un champ de vitesse bidimensionnel $v=(v_x,v_y)$.

\subsection{Imposition faible des contraintes : pénalisation de la divergence}

Le réseau produit directement $v_\theta(x,y)$, avec une fonction de coût
\[\boxed{\mathcal L_{\rm weak}=\frac1N\sum_{i=1}^N\|v_\theta(x_i)-v_i\|^2+\lambda_{\rm div}\frac1N\sum_{i=1}^N|\partial_x v_{\theta,x}(x_i)+\partial_y v_{\theta,y}(x_i)|^2.}\]
Le gradient automatique doit être activé sur les coordonnées. La divergence n'est pas identiquement nulle : elle est seulement pénalisée, avec un compromis contrôlé par $\lambda_{\rm div}$.

\subsection{Imposition forte des contraintes : formulation par fonction de courant}

Un autre réseau apprend un potentiel scalaire $\psi_\theta(x,y)$, puis construit
\[\boxed{v_\theta=(\partial_y\psi_\theta,-\partial_x\psi_\theta).}\]
Si $\psi_\theta$ est suffisamment régulière,
\[\nabla\cdot v_\theta=\partial_x\partial_y\psi_\theta-\partial_y\partial_x\psi_\theta=0.\]
Cette identité est structurelle, indépendamment de la pondération d'une pénalité; des écarts numériques peuvent subsister. Les dérivées automatiques nécessaires au calcul de $v$ ne doivent pas être désactivées pendant l'inférence.
\begin{center}\begin{tikzpicture}[node distance=12mm,>=Latex,box/.style={draw=navy,rounded corners,align=center,minimum height=9mm},every node/.style={font=\small}]
\node[box] (xy) {$(x,y)$};\node[box,right=of xy] (nn) {Réseau\\$\psi_\theta$};\node[box,right=of nn] (grad) {Autograd\\$(\partial_y\psi,-\partial_x\psi)$};\node[box,right=of grad] (v) {$v_\theta$\\$\nabla\cdot v_\theta=0$};
\draw[->,thick] (xy)--(nn);\draw[->,thick] (nn)--(grad);\draw[->,thick] (grad)--(v);
\end{tikzpicture}\end{center}

\section{Solutions commentées des travaux pratiques}

Les réponses ci-dessous suivent les dix exercices des notebooks du répertoire
\texttt{student/}. Les méthodes et fonctions correspondantes figurent dans
\texttt{reference/}. Les observations numériques précises dépendent de
l'entraînement et de l'exécution des notebooks.

\subsection{Exercice 1 --- Propagation avant du MLP}

\textbf{Question.} Reconstituer la méthode \texttt{forward}, comparer les
biais inductifs du MLP et du CNN et discuter la généralisation à une autre
position de la source.
\\

\textbf{Solution commentée.} Dans \texttt{SpatialMLP}, l'entrée
\texttt{x\_onehot} est de dimension $(B,N)$, où $B$ est la taille du
mini-lot et $N$ le nombre de points (64 par défaut). Chaque point reçoit
deux caractéristiques : sa position $x_i$ et la valeur de l'indicateur
de source à ce point. La concaténation produit $(B,N,2)$, puis
l'aplatissement $(BN,2)$. Le réseau entièrement connecté applique
$2\to64\to64\to1$ et la sortie est remise sous la forme $(B,N)$.
Ce MLP agit point par point : il ne combine pas directement les valeurs
des points voisins dans sa propagation avant. Le CNN applique au contraire
des filtres locaux partagés le long de la grille; ce partage des poids
introduit un biais inductif de localité et de translation (modulo les
effets de bord). Une bonne prédiction pour une position de source
observée à l'entraînement ne garantit donc pas une bonne prédiction
pour une nouvelle position : celle-ci doit être évaluée sur les
exemples de test.

\subsection{Exercice 2 --- Construction du graphe kNN}

\textbf{Question.} Construire \texttt{create\_knn\_graph} avec des
connexions orientées symétriques et expliquer le voisinage accessible.
\\

\textbf{Solution commentée.} Pour chaque nœud $i$, on calcule les
distances aux autres nœuds, on exclut $i$, puis on sélectionne ses
$k$ plus proches voisins. Pour chaque voisin $j$, on ajoute les deux
arêtes orientées $(i,j)$ et $(j,i)$ dans \texttt{edge\_index}.
Ces orientations distinctes ne sont pas des doublons. En revanche,
si $i$ et $j$ se sélectionnent mutuellement, la même arête orientée
peut apparaître plusieurs fois. Avec l'agrégation \texttt{mean}
d'\texttt{EdgeConv}, une occurrence répétée est comptée plusieurs
fois et peut modifier la moyenne. Une couche de propagation permet
de recevoir des messages des voisins immédiats, deux couches des
nœuds à distance au plus deux, etc., sans garantir que toutes les
informations soient effectivement conservées.

\subsection{Exercice 3 --- Message EdgeConv}

\textbf{Question.} Construire le message à partir des caractéristiques
des nœuds source et destinataire et interpréter la profondeur du GNN.
\\

\textbf{Solution commentée.} La méthode \texttt{message(x\_i,x\_j)}
concatène les vecteurs du destinataire $i$ et de la source $j$ :
\[
m_{ij}=\operatorname{MLP}_{\theta}([h_i,h_j]),\qquad
h_i'=\frac{1}{|\mathcal N(i)|}\sum_{j\in\mathcal N(i)}m_{ij}.
\]
Ici $h_i,h_j$ sont les vecteurs de caractéristiques, $\theta$ les
paramètres du MLP, $m_{ij}$ le message adressé à $i$, et
$\mathcal N(i)$ ses voisins (sans répétition dans cette écriture).
Trois couches de propagation permettent une dépendance structurelle
aux informations situées à au plus trois arêtes dans le graphe
utilisé. Elles ne permettent pas, par simple propagation locale,
d'accéder à une caractéristique située à une distance arbitraire.
Cette limite concerne la propagation sur le graphe et non une
information globale éventuellement déjà encodée dans les
caractéristiques d'entrée.

\subsection{Exercice 4 --- Données résidu--erreur}

\textbf{Question.} Construire les couples $(r,e)$, démontrer $Ae=r$
et expliquer la normalisation commune.
\\

\textbf{Solution commentée.} Pour $Au^*=f$, une approximation
$u$ obtenue après un nombre aléatoire de pas de Jacobi (de 0 à 59
dans le code) donne
\[
r=f-Au,\qquad e=u^*-u,\qquad
Ae=Au^*-Au=f-Au=r.
\]
La solution $u^*$ est calculée par solveur direct creux. Le code
conserve les points intérieurs et utilise
$s_r=\sqrt{\operatorname{mean}(r^2)}+10^{-10}$, puis fournit au
CNN les couples $(r/s_r,e/s_r)$. Employer le même facteur préserve
la relation d'échelle entre résidu et erreur.

\subsection{Exercice 5 --- Correction neuronale dans Jacobi}

\textbf{Question.} Compléter \texttt{apply\_cnn\_correction} et
\texttt{run\_hybrid}; expliquer l'ajout de la correction et les
courbes d'erreur.
\\

\textbf{Solution commentée.} Le réseau prédit une approximation
$\widehat e$ de l'erreur à partir du résidu courant, après
normalisation. On dénormalise puis on effectue
$u^+=u+\widehat e$ sur les points intérieurs, en conservant les
conditions au bord. La sortie du CNN est une \emph{correction},
pas une solution complète. Dans le notebook, elle est appliquée
une seule fois au pas \texttt{cnn\_at\_step}, puis Jacobi se poursuit.
La courbe d'erreur relative peut montrer une baisse brusque au
moment de la correction; son efficacité dépend de la qualité de
$\widehat e$ et n'est pas garantie pour tout problème.

\subsection{Exercice 6 --- Données second membre--solution}

\textbf{Question.} Construire les couples $(f,u^*)$ normalisés et
comparer les deux applications apprises.
\\

\textbf{Solution commentée} Dans \texttt{hybrid\_accelerator},
le CNN apprend directement
\[
G_\theta:\ f\longmapsto u^*\approx A^{-1}f.
\]
Le jeu de données contient les valeurs intérieures du second
membre et de la solution de référence, normalisées par
$s_f=\sqrt{\operatorname{mean}(f_{\mathrm{int}}^2)}+10^{-10}$.
Dans \texttt{hybrid\_corrector}, le CNN apprend plutôt
\[
C_\theta:\ r\longmapsto e\approx A^{-1}r,
\]
où $r=f-Au$ dépend de l'itéré courant. Bien que les deux
applications fassent intervenir le même opérateur inverse au
niveau mathématique, elles n'ont ni les mêmes entrées ni le même
rôle algorithmique : initialisation d'un côté, correction d'un
itéré de l'autre.

\subsection{Exercice 7 --- Initialisation neuronale}

\textbf{Question.} Compléter \texttt{apply\_cnn\_initial\_guess}
et la boucle de Jacobi; expliquer les erreurs spectrales persistantes.
\\

\textbf{Solution commentée} Le CNN prédit les valeurs
intérieures de $u^0=G_\theta(f)$; le code fixe les valeurs au
bord à zéro, enregistre l'erreur initiale, puis effectue les
itérations de Jacobi. Pour l'erreur $e^k=u^*-u^k$,
\[
e^{k+1}=(I-D^{-1}A)e^k.
\]
Dans les problèmes elliptiques discrétisés usuels, Jacobi
atténue souvent rapidement les composantes oscillantes
(hautes fréquences), mais plus lentement les composantes
lisses (basses fréquences). Une prédiction initiale de bonne
qualité en norme globale peut encore contenir des composantes
lentes à éliminer. Le test à fréquence élevée sans
réentraînement permet également de discuter la généralisation
hors de la distribution d'entraînement.

\subsection{Exercice 8 --- Incompressibilité faible}

\textbf{Question.} Implémenter la perte sur les données et la
pénalisation de divergence.
\\

\textbf{Solution commentée.} Pour une vitesse prédite
$v_\theta=(v_x,v_y)$,
\[
\mathcal L_{\mathrm{data}}
 =\operatorname{mean}\|v_\theta-v_{\mathrm{true}}\|^2,\qquad
\mathcal L_{\mathrm{div}}
 =\operatorname{mean}\left(
 \frac{\partial v_x}{\partial x}
 +\frac{\partial v_y}{\partial y}\right)^2,
\]
\[
\mathcal L=\mathcal L_{\mathrm{data}}
+\lambda_{\mathrm{div}}\mathcal L_{\mathrm{div}}.
\]
Le coefficient $\lambda_{\mathrm{div}}\geq0$ règle le compromis
entre fidélité aux données et pénalisation de la divergence.
La divergence est calculée par différentiation automatique
par rapport aux coordonnées spatiales; pendant l'entraînement,
\texttt{create\_graph=True} permet de différencier ensuite la
perte par rapport aux paramètres du réseau.

\subsection{Exercice 9 --- Incompressibilité forte}

\textbf{Question.} Reconstruire la vitesse à partir d'une fonction
de courant et démontrer l'annulation de la divergence.
\\

\textbf{Solution commentée.} Si le réseau prédit une fonction
de courant $\psi_\theta(x,y)$, on définit
\[
v_x=\frac{\partial\psi_\theta}{\partial y},
\qquad
v_y=-\frac{\partial\psi_\theta}{\partial x}.
\]
Pour une fonction suffisamment régulière ($C^2$),
\[
\nabla\cdot v_\theta
=\frac{\partial^2\psi_\theta}{\partial x\,\partial y}
-\frac{\partial^2\psi_\theta}{\partial y\,\partial x}=0.
\]
L'incompressibilité est ainsi imposée par la paramétrisation
elle-même, sous réserve des erreurs numériques et de
différentiation automatique.

\subsection{Exercice 10 --- Entraînement et validation}

\textbf{Question.} Compléter les branches d'entraînement faible
et fort; expliquer la nécessité des dérivées spatiales.
\\

\textbf{Solution commentée.} Dans la branche faible,
\texttt{compute\_loss} calcule la perte sur les données et
la pénalisation de divergence. Dans la branche forte, le
réseau prédit $\psi_\theta$ et la vitesse est construite par
dérivation spatiale avant de calculer la perte sur les vitesses.
Pour mettre à jour $\theta$, la rétropropagation doit
traverser cette dérivation : les opérations spatiales doivent
donc conserver un graphe différentiable, notamment grâce à
\texttt{create\_graph=True} pendant l'entraînement.
La validation évalue le modèle sans mise à jour des
paramètres; elle ne nécessite pas de conserver le graphe
de dérivation par rapport aux paramètres, mais la
construction de la vitesse forte exige toujours des
dérivées par rapport aux coordonnées. Il ne faut donc pas
supprimer indistinctement toute différentiation automatique
dans cette branche de validation.

\end{document}
